% Nouveau EGA V (cote 25): archive cote 25, site page 101 (PDF page 102).
% typescript. Draft compared: none.
% AI-transcribed by Claude Opus 5.5 (claude-opus-5-5, adaptive thinking, effort high), Message Batches API, batch ega25b (msgbatch_01A7n5YuXubUPNK1zsepi8RG), 2026-11-13; not yet checked by eye.
\scanpage{25/101}
\mspage{4}

% NOTE: long hand-drawn diagonal strokes run across the page: one in the upper left margin (by the first lines), one from the top right down through lines 1--6, a short one through line 8, and a long one from ``Variante'' down to the foot of the page; they look like strokes marking the passages as dealt with, not deletions of single words
est isomorphe à un $\underline{O}_P(n)$. \ -- \ Du point de vue rédaction, il était mala-%
droit de commencer par considérer un $\underline{L}$ de type fini. Commencer carrément
par considérer \ $L = \Gamma(E^{r+1}-(0),\ q_*(\underline{L}))$ \ définissant le Module
$\tilde{L} \ = \unsure{q}_*(q^*(\underline{L}))$ \ , on sait (Chap II) que \ $L$ estb un Module gradué qui
détermine précisément $\underline{L}$, puis on montre que $L$ est libre\inserted{de rang 1} \deleted{comme Module}
gradué \ par le raisonnement indiqué\hand{:}

b) \ Cas général. \ Il se déduit du cas a) \ grâce à \ III 4.6.5. , en
utilisant la relation \ $H^1(P^r_k,\underline{O}_{P^r_k}) = 0$ \ établie dans III 2 . \ \ CQFD

Variante à \ 1.4. \ Soit \ $\underline{Z}\,(S)$ \ l'ensemble des \ fonctions localement
constantes à valeurs entières sur S, on définit un homomorphisme évident
\[ Z(S) \longrightarrow \mathrm{Pic}(P) \ , \]
\inserted{les} \xout{xxxxx} $\underline{n} \in Z(S)$ \ correspondant en eff\xout{etx} \inserted{aux} \xout{xxx} partitions $(S_n)_{n\in\underline{Z}}$ \ en ouvert\unsure{s}
disjoints (dont certains peuvent être vides), \xout{xxxfaçonxxque} et \xout{xx} à une
telle partition on associe le Module inversible $\underline{O}_P(\underline{n})$ dont la \inserted{restr à} \xout{valeurxsur}
$f^{-1}(S_n)$ \ est $\underline{O}_P(n)$. \ On trouve donc une variante
(x bis) \hfill $\mathrm{Pic}(S)\mathrm{x}\underline{Z}(S) \longrightarrow \mathrm{Pic}(P)$ \ , \hfill
et \ un énoncé plus général et plus satisfaisant que 1.4. affirme que
(sous les conditions de 1.1.) c'est là une bijection. (NB Lorsque S $\neq \emptyset$
\inserted[typed]{resp. connexe}\hand{)}
alors l'application \ canonique $\underline{Z} \to \underline{Z}(S)$ \ associant à tout $n \in \underline{Z}$ \ la foncti-%
on constante de valeur n, est injective resp. surjective, \ et on retrouve .
formellement 1.4. , qu'il y a donc lieu d'énoncer sous la forme plus
générale que je viens d'expliciter . \ Signaler aussi en remarque que
dans le langage du schéma de Picard qui sera introduit au Chap V, 1.1.
(sous la forme équivalente \ précédente) s'énonce simplement \ en disant que
l'homomorphisme canonique
\[ \underline{Z}_S \longrightarrow \underline{\mathrm{Pic}}_{P/S} \]
du schéma en groupes constant \xout{xxxxZxSxdéfinixpar} $\underline{Z}$ sur S dans le schéma

\margin{101}
